\chapter{Introduction}\label{ch:introduction}
Bonds finance borrowing through promised payments. Their valuation depends on when those payments occur and how future amounts are discounted. A single discount yield cannot generally represent rates that vary by maturity; a yield curve records rates against maturity under a stated convention. Price levels alone do not establish rate sensitivity.

This dissertation examines transparent fixed-rate valuation, Nelson--Siegel--Svensson modelling and deterministic interest-rate risk, asking:
\begin{enumerate}
\item\label{rq:bond} How do explicit conventions determine bond prices and yield sensitivities?
\item\label{rq:nss} How can a smooth parametric curve be fitted reliably?
\item\label{rq:risk} How can parallel, shaped and key-rate exposures be represented consistently?
\item\label{rq:verification} What independent evidence supports computational correctness?
\end{enumerate}
The scope comprises regular option-free bonds, whose contracts contain no rights to alter payments, and illustrative offline rates. Credit events, forecasting, trading and backtesting are excluded. The contribution is independent verification of established mathematics, not new financial theory.

\Cref{ch:foundations,ch:bonds,ch:nss,ch:risk} develop the financial argument; \cref{ch:computation,ch:verification,ch:results} present computation and evidence. \Cref{ch:limitations,ch:conclusion} delimit the conclusions; \cref{app:derivations,app:reproducibility} supply derivations and reproduction instructions.
